\scanpage{65}
% @ top, boxed paragraph | text
\fbox{\parbox{.91\linewidth}{La CSG sera donc réincorporée dans le montant net à déclarer qui figurera sur la déclaration fiscale que nous vous adresserons début 1992 pour vos revenus de 1991.}}

% @ upper third, left label and heading | text
\margin{Les personnes exonérées de la CSG}
Sont exonérés de la CSG :

% @ upper third, first dash | text
- les retraités fiscalement non domiciliés en France

% @ upper third, indented paragraph | text
\textbf{A titre transitoire}, il a été admis que les personnes domiciliées à l'étranger seraient réputées être dans cette situation (sans avoir dans l'immédiat à fournir d'autre preuve).

% @ upper third, second dash | text
- Les retraités :

% @ middle, first dot | text
. dont la cotisation d'impôt sur les revenus de l'année précédente n'est pas mise en recouvrement,

% @ middle, second dot | text
. titulaires d'un avantage vieillesse non contributif attribué sous conditions de ressources.

% @ middle, paragraph after list | text
Ces cas d'exonération sont donc identiques à ceux retenus pour la cotisation \textquotedbl{}\textbf{assurance-maladie}\textquotedbl{} de \textbf{2,40 \% cotisation qui reste toujours applicable}, mais les périodes d'exonération sont différentes.

% @ middle, first Si paragraph | text
Si vous êtes actuellement exonéré(e) de la retenue assurance maladie, vous le serez également pour la CSG.

% @ middle, second Si paragraph | text
Si vous êtes devenu(e) non imposable en 1990 au titre des revenus de 1989 alors que vous l'étiez au titre de l'année précédente et :

% @ lower half, first dash | text
- si vous avez justifié cette situation auprès de la Caisse,(a):

% @ lower half, first dot | text
. vous serez exonéré(e) de la CSG dès le paiement de votre retraite du 1er trimestre 1991,

% @ lower half, second dot | text
. vous ne serez exonéré(e) de la retenue assurance maladie qu'à compter du paiement du 3ème trimestre 1991.

% @ lower third, second dash | text
- si vous n'avez pas justifié encore cette situation auprès de la Caisse, nous vous invitons à le faire pour vous permettre de bénéficier de ces exonérations.

% @ lower third, boxed reminder | text
\fbox{\parbox{.91\linewidth}{Nous vous rappelons que lorsque vous redevenez à nouveau imposable, vous devez en informer la Caisse.}}

% @ foot, note (a) | text
(a) par l'envoi de votre avis de non imposition
% archive sheet number: 64
